\begin{textsample}{POS Dim 1 – ac – Score 36.00 – ac/ac\_em\_8.txt}  \label{ex:f1_pos_195}
7.
Protagonismo juvenil A escola, imersa na sociedade \textbf{contemporânea}, recebe estudantes com comportamentos pós-modernos, isso \textbf{implica} em \textbf{dizer}, \textbf{que} o ensino oferecido a esses jovens deve proporcionar mecanismos \textbf{que} estimulem e contemplem o desenvolvimento da autonomia intelectual, do pensamento crítico e reflexivo, bem como o aprimoramento do \textbf{educando} como pessoa humana, levando em consideração a sua formação ética, e seu protagonismo, em prol de \textbf{uma} sociedade mais \textbf{justa}, ética, democrática, inclusiva, sustentável e solidária ( BNCC, 2018 ), princípios essenciais para a participação cidadã dos jovens e sua formação profissional, visando atender às exigências da \textbf{contemporaneidade}.
Dentro do contexto pós \textbf{moderno}, tendo como foco a formação dos estudantes no ensino médio, as DCNEM/2011 \textbf{explicitam} e acentuam sobre a importância de \textbf{uma} não caracterização homogênea do público pertencente a essa etapa, público esse \textbf{que} é composto por jovens e adolescentes diversos, sobre os quais, não se deve conceber e reproduzir a ideia de “ juventude ” como \textbf{um} momento transitório de passagem, entre a infância e a \textbf{maturidade}, mas sim como \textbf{uma} condição sócio-histórico-cultural de \textbf{uma} \textbf{categoria} de \textbf{sujeitos} \textbf{que} \textbf{necessita} ser considerada em suas múltiplas dimensões, com especificidades próprias \textbf{que} não estão restritas às dimensões biológica e etária, mas \textbf{que} se encontram articuladas com \textbf{uma} multiplicidade de atravessamentos sociais e culturais, produzindo múltiplas culturas juvenis ou muitas juventudes. ( Parecer CNE/CEB nº 5/2011 ).
Reconhecer as juventudes, significa entendê -las como diversificadas e plurais, considerando \textbf{que} os agentes participantes delas, são ativos na sociedade em \textbf{que} estão inseridos, são sujeitos sociais participativos, e \textbf{que} essa participação direcionada e colaborativa com os adultos e educadores pode gerar mudanças na realidade individual, social, ambiental, cultural e política, acentuando valores de autonomia e autodeterminação.
A concepção de educação apresentada pela ONU, no início do \textbf{século} XXI, tem como \textbf{fundamento} o paradigma do desenvolvimento humano, \textbf{que} vem sendo difundido desde 1990 pelo PNUD ( Programa das Nações Unidas para o Desenvolvimento ).
Segundo essa visão, a escola tem como objetivo criar espaços e condições para \textbf{que} os jovens possam envolver -se em atividades \textbf{que} sejam direcionadas para a resolução de problemas do cotidiano, passando a atuar como protagonistas, a partir de suas próprias \textbf{percepções} da realidade em \textbf{que} se inserem.
Nesse sentido, o protagonismo juvenil, \textbf{enquanto} modalidade de ação educativa, é a criação de espaços e condições capazes de possibilitar aos jovens envolver -se em atividades direcionadas à solução de problemas reais, atuando como fonte de iniciativa, liberdade e \textbf{compromisso}. ( Costa, 2006 ).
A palavra “ protagonista ”, usada comumente na literatura e no teatro, recentemente foi ressignificada na sociologia e na política para acentuar os agentes sociais como personagens principais de vários movimentos e atuações.
Na educação, o termo é relacionado à participação ativa do estudante no seu processo educacional como agente \textbf{corresponsável}, o \textbf{que} \textbf{exige} ações pedagógicas constantes \textbf{que} permitam ao estudante o desenvolvimento de competências e habilidades \textbf{que} visem \textbf{uma} aprendizagem além da “ \textbf{memorização} ” dos conteúdos e o treinamento para respostas certas, voltada para a formação de \textbf{sujeitos} ativos, capazes de tomar decisões e fazer escolhas \textbf{embasadas} no conhecimento, na reflexão, levando em consideração a si próprio e o coletivo envolto.
Dessa forma, “ os educadores passam a chamar de protagonismo os processos, movimentos e dinamismos sociais e educativos, nos quais os adolescentes e jovens assumem o papel principal das ações \textbf{que} executam ” ( ISE, 2019 ).
No Novo Ensino Médio, a introdução do protagonismo como \textbf{um} princípio educativo, tendo o estudante a \textbf{centralidade} na elaboração do seu projeto de vida, está \textbf{ancorado} nos quatros pilares da educação, apresentados pela Comissão Internacional sobre a Educação para o Século XXI, da Unesco, no início do década de 1990, onde se pensavam perspectivas, concepções e práticas pedagógicas frente às mudanças do novo \textbf{século}, com a intensificação da oferta e dos meios para circulação do conhecimento, armazenamento de informações e comunicações.
No relatório da comissão, a Unesco convida a pensar a educação ao longo da vida como \textbf{uma} maneira de lidar e \textbf{viver} no mundo \textbf{contemporâneo}, com suas transformações, ressaltando \textbf{que} a “ educação não se \textbf{restringe} à transmissão de conhecimentos, mas também à criação de \textbf{um} desejo de continuar aprendendo durante toda a vida, a partir do reconhecimento de \textbf{que} aprender é \textbf{viver} em transformações de si próprio e dos outros ” ( UNESCO, 1996 ).
Saber “ conhecer, fazer, ser e conviver ” são aprendizagens \textbf{que} se constroem através de situações pedagógicas, \textbf{que} \textbf{explícita} ou implicitamente têm o protagonismo como finalidade.
O protagonismo, como princípio educativo no Novo Ensino Médio, faz parte de \textbf{uma} concepção cidadã, de escola para as juventudes, assim como os princípios da \textbf{contextualização}, da trans/interdisciplinaridade e da aprendizagem colaborativa, \textbf{que} estão contidos no parágrafo 2º, do artigo 7º das Diretrizes Curriculares Nacionais para o Ensino Médio, em \textbf{que} se \textbf{explicita} \textbf{que} : > (... ) o currículo deve contemplar tratamento \textbf{metodológico} \textbf{que} evidencie a \textbf{contextualização}, a diversificação e a \textbf{transdisciplinaridade} ou formas de interação e articulação entre diferentes campos de saberes específicos, contemplando vivências e práticas e vinculado à educação escolar ao mundo do trabalho e à prática social. ( Parecer CNE/CEB nº 5/2011 ).
Os jovens são pessoas em pleno desenvolvimento, \textbf{que} \textbf{necessitam} conhecer e experimentar a partir de suas próprias convicções, para compreenderem o mundo e a sociedade da qual fazem parte.
A possibilidade de vivenciarem escolhas, refletirem se foram adequadas ou não aos contextos, favorece o processo de autorregulação, o “ aprender a aprender ”.
Erik Erikson, psicanalista alemão, em seu estudo sobre a adolescência, a partir de sua \textbf{teoria} do Desenvolvimento Psicossocial, \textbf{argumenta} \textbf{que} à medida \textbf{que} o jovem amplia sua convivência e participação social, inicia a experimentação de novos papéis, cria, estabelece laços e vínculos tanto na esfera pessoal quanto na dimensão produtiva e social, o jovem começa a ampliar sua rede de contatos e a perceber \textbf{que} \textbf{uma} nova maneira de se relacionar com as demandas passa a ser \textbf{exigida}, porque para cada \textbf{um} dos novos segmentos \textbf{que} passa a fazer parte de sua vida, há \textbf{um} impacto próprio exercido por ele.
É justamente nesse momento \textbf{que} \textbf{um} senso de coerência entre aquilo \textbf{que} acredita, deseja, faz e planeja se faz importante.
A construção do seu próprio repertório de valores, critérios e crenças morais compõe \textbf{uma} espécie de “ ideologia pessoal ” ( apud ICE, 2019 ).
Desta forma, promover a escuta dos estudantes é \textbf{um} passo fundamental para apoiar a tomada de decisões, \textbf{que} vão desde o planejamento de atividades e da adoção de novas práticas pedagógicas até \textbf{tópicos} mais complexos, como a parte flexível deste currículo escolar.
O estudante poderá escolher, conforme seu interesse e aptidões, \textbf{um} percurso formativo \textbf{que} melhor condiz com o seu projeto de vida, tendo a possibilidade de consolidação por meio de orientações, reflexões e aprofundamento de sua aprendizagem ; \textbf{logo}, é importante \textbf{que} a equipe escolar o incentive a refletir sobre quem ele é e como se \textbf{vê} no futuro.
O processo de escuta e protagonismo possibilita ao estudante o exercício de práticas e vivências, nas quais, a partir do conhecimento adquirido, terá condições de analisar fenômenos diferentes e lançar mão de ações propositivas em seu contexto escolar e na sociedade, de modo geral.
Além disso, desenvolverá também condições essenciais para seu crescimento pessoal e coletivo, aspectos necessários para a construção de sua identidade e para o desenvolvimento de sua autoestima, em vários aspectos.
Nesse sentido, o protagonismo potencializa atitudes alicerçadas no conhecimento adquirido e na inovação, para \textbf{que}, a partir do reconhecimento de suas potencialidades, os jovens identifiquem perspectivas possíveis para descobrir novas formas de atuação, ressignificando questões antigas e descobrindo novas formas de pensar e agir diante de situações, cenários e contextos jamais imaginados por eles, sejam no presente ou no futuro, fazendo uso de seu potencial criativo e de sua força transformadora.

% matched lemmas: ancorar, argumentar, categoria, centralidade, compromisso, contemporaneidade, contemporâneo, contextualização, corresponsável, dizer, educar, embasar, enquanto, exigir, explicitar, explícito, fundamento, implicar, justo, logo, maturidade, memorização, metodológico, moderno, necessitar, percepção, que, restringir, sujeito, século, teoria, transdisciplinaridade, tópico, um, ver, viver
\end{textsample}
